% ─────────────────────────────────────────────────────────────────
%  Draft of §IV.A — to replace existing §IV.A (Phase diagram, lines
%  362–429) and §IV.B (Spatiotemporal structure, lines 431–485).
%  Hot-runaway attractor is described later in §IV.D, not here.
% ─────────────────────────────────────────────────────────────────

\section{Results and Discussion}
\label{sec:results}

The dispersion analysis of Sec.~\ref{sec:lsa} establishes the linear
bifurcation conditions for the homogeneous swollen state.  We now
turn to direct simulations of the full PDE system
\eqref{eq:governing} in the planar slab and report four results.
First, default-IC simulations across the principal parameter planes
realize three distinct attractors—cold-swollen, an LCST-front
oscillation, and a frozen partial-collapse front
(Sec.~\ref{sec:phasediagram}).  Second, the LCST-front attractor is
a relaxation oscillator on a slow manifold generated by the surface
boundary condition together with the $\varphi$-dependent diffusion
barrier; the same barrier prevents the gel from undergoing global
collapse (Sec.~\ref{sec:lcst_front_mechanism}).  Third, the 0D
linear stability analysis is a necessary but not sufficient
predictor of the PDE oscillation: a closed-form parameter
inversion of the homogeneous-SS conditions identifies a Whitney-
cusp catastrophe whose collapsed-only branch covers
$\sim\!51\%$ of the displayed $(\Bi_T,S_\chi)$ plane, and the
saddle--node-on-invariant-circle scaling at the upper-$S_\chi$
boundary confirms that the cycle is destroyed by a global
topological event rather than by a local Hopf instability of any
homogeneous branch (Sec.~\ref{sec:lsa_vs_nl}).  Fourth,
initial-condition sweeps reveal a fourth attractor—a hot-runaway
super-swollen steady state—coexisting with the LCST-front cycle in
a substantial portion of the C-region; the bistability gives a
falsifiable IC-induced hysteresis prediction
(Sec.~\ref{sec:multistability}).

\subsection{Realized attractors and their spatiotemporal structure}
\label{sec:phasediagram}

We integrate the PDE system from the canonical cold-swollen initial
condition ($J_0=1.30$, $\theta_0=0$, $u_0=0.02$) to dimensionless
time $t=200$, classify each parameter point by combining the surface
swelling amplitude, the peak polymer fraction, and the oscillation
period into a six-class regime label, and sweep $25\times25$ points
in the $(\Bi_T, S_\chi)$ plane and $25\times15$ points in the
$(\Bi_T, \Da)$ plane (Fig.~\ref{fig:fig4_phase}).

\paragraph{Realized attractors (default-IC).}
Three of the six classification labels are populated:
\begin{itemize}
\item[(i)] a uniform \emph{cold-swollen} steady state at the cold
  edge of the diagram (low $S_\chi$ or high $\Bi_T$);
\item[(ii)] a large-amplitude \emph{LCST-front} oscillation occupying
  the central $\sim\!45\%$ of the $(\Bi_T,S_\chi)$ plane, in which a
  narrow surface shell cycles between collapsed and swollen states
  while the centre stays cold-swollen;
\item[(iii)] a \emph{frozen partial-collapse} state at large $\Bi_T$,
  with the same two-zone spatial structure as the LCST-front cycle
  but stationary in time.
\end{itemize}
The remaining classifier-allowed regimes (globally collapsed steady,
globally collapsing oscillation) never appear in the $1{,}000$
default-IC simulations covered by the two parameter sweeps; the
mechanism that suppresses them is discussed in
Sec.~\ref{sec:lcst_front_mechanism}.  A fourth attractor, a uniform
\emph{hot-swollen runaway} steady state, exists in a substantial
sub-region of the LCST-front map but is inaccessible from the
default cold IC and is revealed only by initial-condition variation
(Sec.~\ref{sec:multistability}).

\paragraph{Spatiotemporal structure of the three default-IC
attractors.}
Figure~\ref{fig:fig_all_regimes} shows the spatiotemporal fields
$J(\xi,t)$, $\theta(\xi,t)$ and the surface and centre $J(t)$ traces
for one representative point per attractor.

The cold-swollen state [Fig.~\ref{fig:fig_all_regimes}(a)] is
spatially uniform with $J\!\simeq\!1.30$ and $\theta\!\simeq\!0$
everywhere; the surface BC is passively satisfied without driving
significant transport.

The LCST-front cycle [Fig.~\ref{fig:fig_all_regimes}(b)] has a
two-zone structure.  The outer shell ($\xi\!\in\![0.885,\,1]$ at
the working point) cycles between the swollen state
($J\!\sim\!1.4$) and the collapsed state ($J\!\sim\!0.15$) with
period $T\!\simeq\!8\text{--}11$.  The inner core ($\xi<0.885$)
stays at $J\!\simeq\!1.39$ throughout the cycle.  The temperature
field tracks the asymmetry: $\theta_\mathrm{surf}$ swings between
$\sim\!1.5$ and $\sim\!4.8$ each cycle, while
$\theta_\mathrm{centre}$ is held nearly constant at
$\theta_\mathrm{centre}\!\simeq\!3.2$ by heat conduction from the
warm shell.

The frozen-front state [Fig.~\ref{fig:fig_all_regimes}(c)] inherits
the same two-zone spatial structure, now stationary in time
($J_\mathrm{surf}\!\simeq\!0.15$,
$J_\mathrm{centre}\!\simeq\!1.4$).  At sufficiently large $\Bi_T$,
the surface heat flux $\Bi_T\theta_\mathrm{surf}$ is large enough
to clamp $\theta_\mathrm{surf}$ below the ignition threshold and
the cycle of Sec.~\ref{sec:lcst_front_mechanism} is locked at
phase~3.
